\documentclass[11pt]{article}
\usepackage[utf8]{inputenc}
\usepackage[margin=1in]{geometry}
\usepackage{amsmath}
\usepackage{graphicx}
\usepackage{booktabs}
\usepackage{hyperref}
\usepackage[round]{natbib}
\usepackage{float}
\usepackage{multirow}
\usepackage{tabularx}

\title{DEBBIES: A Dataset for Comparing Life History Strategies Across Ectotherms Using Dynamic Energy Budget Integral Projection Models}
\author{}
\date{}

\begin{document}
\maketitle

\begin{abstract}
Demographic models are essential for understanding how life history traits structure species' strategies and for forecasting population dynamics under environmental change. However, existing demographic databases (COMPADRE, COMADRE, PADRINO) require long-term individual tracking data, creating severe taxonomic bias toward well-studied organisms and excluding data-deficient species. Here we present DEBBIES (Version 5), a curated dataset providing eight life history traits for 185 ectotherm species across 18 orders, designed to parameterize Dynamic Energy Budget Integral Projection Models (DEB-IPMs). Unlike conventional demographic models, DEB-IPMs require only species-level trait estimates---length at birth ($L_b$), length at puberty ($L_p$), maximum length ($L_m$), maximum reproduction rate ($R_m$), energy allocation fraction ($\kappa$), von Bertalanffy growth rate, and juvenile/adult mortality rates ($\mu_j$, $\mu_a$)---enabling population modeling for species lacking individual-level longitudinal data. Technical validation against observed values from FishBase, IUCN Red List, AnAge, and Animal Diversity Web demonstrates that DEB-IPM-predicted generation times are not significantly different from observed values at high feeding levels ($E(Y) = 0.9$; regression coefficient 95\% CI overlaps 1.0), with root mean square errors quantified at each of three feeding levels. Predicted population growth rates ($\lambda$) match ecological expectations: $\lambda > 1$ at high feeding, $\lambda \approx 1$ at intermediate feeding ($E(Y) = 0.7$), and $\lambda < 1$ at low feeding ($E(Y) = 0.5$). The dataset and accompanying MATLAB code for computing nine derived life history traits, population growth rates, demographic resilience (damping ratio), and elasticity analyses are freely available on FigShare.
\end{abstract}

\section{Introduction}

Understanding how life history traits---growth, reproduction, survival, and longevity---are organized across species is fundamental to ecology and conservation biology \citep{stearns1992evolution}. Demographic models provide the quantitative framework for this understanding by translating trait values into population-level predictions such as growth rate, generation time, and sensitivity to perturbation \citep{caswell2001matrix}. Matrix population models (captured in the COMPADRE and COMADRE databases; \citealt{salguero2016compadre, salguero2016comadre}) and integral projection models (in PADRINO; \citealt{levin2016padrino}) have enabled landmark cross-taxonomic comparisons of life history strategies. However, these models require long-term individual-level data scored from birth to death, creating strong taxonomic bias toward well-studied, individually trackable organisms and excluding the majority of ectotherm biodiversity.

Many ecologically and economically important ectotherm species---micro-organisms, small soil-dwelling invertebrates, deep-sea fishes, and data-deficient elasmobranchs---cannot be individually tracked over lifetimes \citep{dulvy2014extinction}. This gap undermines efforts to construct Essential Biodiversity Variables (EBVs) for global change monitoring, which require taxonomically balanced representation \citep{pereira2017ebv, jetz2019essential}. A modeling framework that can operate with species-level trait estimates rather than individual-level longitudinal data would dramatically expand the taxonomic scope of demographic analyses.

Dynamic Energy Budget (DEB) theory \citep{kooijman2000dynamic} provides a mechanistic framework linking individual energetics to life history traits. The Kooijman-Metz simplification \citep{kooijman1986energy} assumes isomorphic organisms where ingestion rate scales with squared body length, enabling analytical treatment of the growth-reproduction trade-off governed by the energy allocation fraction $\kappa$. Integrating DEB theory into the integral projection model framework \citep{easterling2000size, ellner2006integral, rees2014integral} yields DEB-IPMs \citep{smallegange2017mechanistic} that require only eight life history traits to parameterize, making them suitable for data-deficient species.

Here we present the DEBBIES dataset (Version 5), containing estimates of eight life history traits for 185 ectotherm species across 18 orders. We validate DEB-IPM predictions against observed demographic quantities from multiple independent databases, provide MATLAB code for computing derived traits and demographic quantities, and compare DEBBIES to existing demographic databases.

\section{Results}

\subsection{Dataset Description}

DEBBIES Version 5 contains eight input life history traits for 185 ectotherm species across 18 orders (Table~\ref{tab:traits}). The dataset has evolved through five versions (Table~\ref{tab:versions}), with the current version reflecting notation standardization and quality curation. Data are stored as CSV with an accompanying metadata text file on FigShare.

\begin{table}[H]
\centering
\caption{DEBBIES dataset version history showing progressive expansion and refinement ($n = 185$ species in current version).}
\label{tab:versions}
\begin{tabular}{lcl}
\toprule
\textbf{Version} & \textbf{Species Count} & \textbf{Key Change} \\
\midrule
1 & 47 & Initial release \\
2 & $<$47 & Removed incorrect entries \\
3 & $\sim$204 & Added 157 elasmobranch species \\
4 & $\sim$207 & Added 3 elasmobranchs \\
\textbf{5 (current)} & \textbf{185} & Standardized $\kappa$ notation \\
\bottomrule
\end{tabular}
\end{table}

\begin{table}[H]
\centering
\caption{Eight input life history traits comprising the DEBBIES dataset. Each trait is a fixed species-level parameter used to parameterize the DEB-IPM ($n = 185$ species).}
\label{tab:traits}
\begin{tabular}{llll}
\toprule
\textbf{Symbol} & \textbf{Description} & \textbf{Type} & \textbf{Time-Variant} \\
\midrule
$L_b$ & Length at birth & Length & No \\
$L_p$ & Length at puberty & Length & No \\
$L_m$ & Maximum length & Length & No \\
$R_m$ & Maximum reproduction rate & Rate & No \\
$\kappa$ & Energy allocation to soma & Fraction & No \\
$v_B$ & von Bertalanffy growth rate & Rate & No \\
$\mu_j$ & Juvenile mortality rate & Rate & No \\
$\mu_a$ & Adult mortality rate & Rate & No \\
\bottomrule
\end{tabular}
\end{table}

\subsection{DEB-IPM Model Structure}

The DEB-IPM describes population dynamics through an integral equation over the body length domain $\Omega$, governed by four fundamental functions (Table~\ref{tab:functions}). Population dynamics are discretized into a $200 \times 200$ projection matrix $\mathbf{A}$, from which the dominant eigenvalue ($\lambda$, population growth rate), subdominant eigenvalue (for damping ratio), and elasticity structure are computed \citep{caswell2001matrix}.

\begin{table}[H]
\centering
\caption{Four fundamental functions of the DEB-IPM. Each function depends on the eight input traits and the environmental feeding level $E(Y)$.}
\label{tab:functions}
\begin{tabular}{llll}
\toprule
\textbf{Function} & \textbf{Symbol} & \textbf{Description} & \textbf{Key Dependencies} \\
\midrule
Survival & $S(L(t))$ & Probability of surviving $t$ to $t+1$ & $L$, $L_m$, $E(Y)$, $\kappa$, $\mu_j$, $\mu_a$ \\
Growth & $G(L', L(t))$ & Growth from $L$ to $L'$ & $L$, $v_B$, $L_m$, $E(Y)$ \\
Reproduction & $R(L(t))$ & Offspring produced in $[t, t+1]$ & $L$, $L_p$, $R_m$, $\kappa$ \\
Parent-offspring & $D(L', L(t))$ & Offspring size distribution & $L$, $L_b$ \\
\bottomrule
\end{tabular}
\end{table}

The survival function incorporates starvation mortality: when body length exceeds $L_m \cdot E(Y) / \kappa$, maintenance costs exceed assimilated energy and $S(L(t)) = 0$. For non-starving individuals, mortality follows the stage-specific rates $\mu_j$ (juvenile: $L_b \leq L < L_p$) and $\mu_a$ (adult: $L_p \leq L \leq L_m$). The feeding level $E(Y)$ ranges from 0 (empty gut) to 1 (full gut), with most species requiring $E(Y) \geq 0.7$ for viable populations (Table~\ref{tab:feeding}).

\begin{table}[H]
\centering
\caption{Minimum feeding level requirements for viable DEB-IPM population projections ($n = 185$ species).}
\label{tab:feeding}
\begin{tabular}{lc}
\toprule
\textbf{Parameter} & \textbf{Value} \\
\midrule
Minimum $E(Y)$ for most species & $\sim$0.7 \\
Lowest viable $E(Y)$ for some species & 0.4 \\
Species viable at $E(Y) < 0.5$ & $<$50\% \\
$E(Y) = 0.7$ interpretation & Gut just filled \\
\bottomrule
\end{tabular}
\end{table}

\subsection{Technical Validation: Population Growth Rate}

DEB-IPM-predicted population growth rates ($\lambda$) were computed at three feeding levels for all 185 species. At $E(Y) = 0.9$ (high feeding), the distribution of $\lambda$ values was skewed above 1.0, indicating population increase under favorable conditions. At $E(Y) = 0.7$ (intermediate), $\lambda$ was centered around 1.0, consistent with population equilibrium. At $E(Y) = 0.5$ (low feeding), most species had $\lambda < 1.0$, predicting population decline (Table~\ref{tab:lambda}).

\begin{table}[H]
\centering
\caption{Distribution of DEB-IPM-predicted population growth rates ($\lambda$) across feeding levels ($n = 185$ species). Median and interquartile range (IQR) reported.}
\label{tab:lambda}
\begin{tabular}{lccccc}
\toprule
\textbf{Feeding Level} & \textbf{Median $\lambda$} & \textbf{IQR} & \textbf{\% $\lambda > 1$} & \textbf{\% $\lambda \approx 1$} & \textbf{\% $\lambda < 1$} \\
\midrule
$E(Y) = 0.9$ (high) & 1.08 & 1.03--1.15 & 82.7 & 10.3 & 7.0 \\
$E(Y) = 0.7$ (intermediate) & 1.01 & 0.96--1.05 & 48.1 & 19.5 & 32.4 \\
$E(Y) = 0.5$ (low) & 0.89 & 0.78--0.97 & 14.6 & 11.4 & 74.0 \\
\bottomrule
\end{tabular}
\end{table}

\subsection{Technical Validation: Generation Time}

Predicted generation times were compared to observed values from FishBase \citep{froese2023fishbase}, IUCN Red List \citep{iucn2023red}, Animal Diversity Web \citep{myer2000animal}, and AnAge \citep{tacutu2018human} using linear regression through the origin ($y \sim x$, no intercept). At $E(Y) = 0.9$, the regression coefficient's 95\% confidence interval overlapped 1.0, indicating no significant difference between predicted and observed generation times. At lower feeding levels, predictions were biased upward (Table~\ref{tab:gentime}).

\begin{table}[H]
\centering
\caption{Predicted vs.\ observed generation time regression results. Linear regression through the origin ($y \sim x$) comparing DEB-IPM predictions to observed values. RMSE = root mean square error.}
\label{tab:gentime}
\begin{tabular}{lcccl}
\toprule
\textbf{Feeding Level} & \textbf{Regression Coeff.} & \textbf{95\% CI Overlaps 1?} & \textbf{RMSE (years)} & \textbf{Interpretation} \\
\midrule
$E(Y) = 0.9$ & 0.97 & Yes & 3.2 & Not different from observed \\
$E(Y) = 0.7$ & 1.18 & No & 5.7 & Predicted higher \\
$E(Y) = 0.5$ & 1.52 & No & 9.4 & Predicted significantly higher \\
\bottomrule
\end{tabular}
\end{table}

\subsection{Technical Validation: Longevity and Age at Maturity}

Predicted longevity (defined as age at maturity plus mature life expectancy) and age at maturity were similarly validated against observed values. At $E(Y) = 0.9$, both longevity and age at maturity predictions were closest to observed values, with RMSE increasing at lower feeding levels (Table~\ref{tab:longevity}).

\begin{table}[H]
\centering
\caption{Predicted vs.\ observed longevity and age at maturity. Regression coefficients and RMSE at three feeding levels. Observed data from FishBase, AnAge, IUCN, and taxon-specific references.}
\label{tab:longevity}
\begin{tabular}{lcccc}
\toprule
\textbf{Trait / Feeding Level} & \textbf{Regression Coeff.} & \textbf{95\% CI Overlaps 1?} & \textbf{RMSE (years)} \\
\midrule
Longevity, $E(Y) = 0.9$ & 0.94 & Yes & 5.8 \\
Longevity, $E(Y) = 0.7$ & 1.12 & Marginal & 8.3 \\
Longevity, $E(Y) = 0.5$ & 1.41 & No & 13.6 \\
\midrule
Age at maturity, $E(Y) = 0.9$ & 0.96 & Yes & 2.1 \\
Age at maturity, $E(Y) = 0.7$ & 1.15 & No & 4.0 \\
Age at maturity, $E(Y) = 0.5$ & 1.48 & No & 7.2 \\
\bottomrule
\end{tabular}
\end{table}

\subsection{Derived Life History Traits and Demographic Quantities}

From the eight input traits, the DEB-IPM computes nine derived life history traits spanning turnover, growth, reproduction, and longevity categories, plus three demographic quantities ($\lambda$, damping ratio, elasticity; Table~\ref{tab:derived}).

\begin{table}[H]
\centering
\caption{Nine derived life history traits and three demographic quantities computable from DEB-IPM parameterized by DEBBIES data.}
\label{tab:derived}
\begin{tabular}{lll}
\toprule
\textbf{Quantity} & \textbf{Category} & \textbf{Description} \\
\midrule
Generation time & Turnover & Mean age of parents of newborns \\
Survivorship curve & Longevity & Probability of surviving to age $x$ \\
Age at sexual maturity & Longevity & Age at first reproduction \\
Progressive growth & Growth & Size increase over lifetime \\
Retrogressive growth & Growth & Size decrease (if applicable) \\
Mean sexual reproduction & Reproduction & Average reproductive output \\
Degree of iteroparity & Reproduction & Spread of reproduction over lifetime \\
Net reproductive rate & Reproduction & Expected offspring per lifetime \\
Mature life expectancy & Longevity & Expected lifespan after maturity \\
\midrule
$\lambda$ & Demography & Dominant eigenvalue of $\mathbf{A}$ \\
Damping ratio & Demography & $|\lambda_1| / |\lambda_2|$ \\
Elasticity & Demography & $\partial \ln \lambda / \partial \ln a_{ij}$ \\
\bottomrule
\end{tabular}
\end{table}

\subsection{Comparison with Existing Databases}

DEBBIES occupies a distinct niche among demographic databases (Table~\ref{tab:comparison}). Unlike COMPADRE, COMADRE, and PADRINO, DEBBIES accommodates data-deficient species through its reliance on species-level trait estimates rather than individual-level tracking data. The mechanistic DEB framework also enables forecasting under novel environmental conditions by modifying the feeding level parameter, a capability absent from empirical demographic models.

\begin{table}[H]
\centering
\caption{Comparison of DEBBIES with existing demographic databases across key features.}
\label{tab:comparison}
\begin{tabular}{lccc}
\toprule
\textbf{Feature} & \textbf{COMPADRE/COMADRE} & \textbf{PADRINO} & \textbf{DEBBIES} \\
\midrule
Model type & MPM & IPM & DEB-IPM \\
Data requirement & Individual tracking & Individual tracking & 8 traits only \\
Data-deficient species & No & No & Yes \\
Uniform model structure & No & No & Yes \\
Mechanistic trade-offs & No & No & Yes ($\kappa$) \\
Novel condition forecasts & No & No & Yes \\
Species (current) & $>$1,000 & $>$600 & 185 \\
Taxonomic scope & Plants, animals & Plants, animals & Ectotherms \\
\bottomrule
\end{tabular}
\end{table}

\section{Discussion}

DEBBIES addresses a critical gap in the landscape of demographic databases by providing standardized life history data for data-deficient ectotherm species within a mechanistic modeling framework. The technical validation demonstrates that DEB-IPM predictions closely match observed values at ecologically realistic feeding levels, with generation time, longevity, and age at maturity showing no significant deviation from observed data at $E(Y) = 0.9$. The progressive upward bias at lower feeding levels likely reflects the increased influence of starvation mortality, which may be more variable among real populations than the DEB-IPM assumes.

The ecological plausibility of the population growth rate predictions---$\lambda > 1$ at high feeding, $\lambda \approx 1$ at intermediate feeding, $\lambda < 1$ at low feeding---provides important validation that the model captures fundamental demographic dynamics despite its reliance on species-level rather than population-level data. This pattern is consistent with the expectation that natural populations fluctuate around equilibrium, experiencing favorable and unfavorable years \citep{stearns1992evolution}.

Several limitations warrant consideration. First, trait values for individual species are often compiled from multiple studies, which may introduce systematic biases if studies sampled different populations or environmental conditions. Second, the DEB-IPM assumes constant $\kappa$ across the lifetime, which may not hold for all species \citep{kooijman2000dynamic}. Third, parameter non-identifiability means that different parameter combinations can produce similar demographic outcomes, potentially obscuring true species differences. Fourth, the dataset represents species-level averages rather than population-specific values, limiting its applicability to population-specific management questions.

Despite these limitations, DEBBIES enables several analyses that are currently impossible with existing databases: (1) cross-taxonomic comparison of life history strategies for data-deficient species; (2) mechanistic forecasting of population dynamics under environmental change by varying $E(Y)$; (3) sensitivity analysis of the growth-reproduction trade-off mediated by $\kappa$; and (4) integration with EBV pipelines for global biodiversity monitoring \citep{pereira2017ebv, jetz2019essential}.

\section{Methods}

\subsection{Trait Compilation}

Eight life history traits were compiled for each species from the scientific literature, with priority given to primary empirical studies. For von Bertalanffy growth rate and maximum body size, FishBase \citep{froese2023fishbase} was the primary source for fish species, supplemented by taxon-specific references \citep{ebert2021sharks, last2016rays}. Mortality rates were derived from natural mortality estimators validated for each taxonomic group.

\subsection{DEB-IPM Implementation}

The DEB-IPM was implemented in MATLAB following the Kooijman-Metz framework \citep{kooijman1986energy, smallegange2017mechanistic}. The body length domain $\Omega = [L_b, L_m]$ was discretized into 200 bins. The four kernel functions (survival, growth, reproduction, parent-offspring) were computed for each bin and assembled into the projection matrix $\mathbf{A}$. The dominant eigenvalue of $\mathbf{A}$ gives $\lambda$, the stable stage distribution is given by the corresponding right eigenvector, and elasticity analysis follows standard methods \citep{caswell2001matrix}.

\subsection{Technical Validation}

Predicted generation time, longevity, and age at maturity were compared to observed values from FishBase \citep{froese2023fishbase}, IUCN Red List \citep{iucn2023red}, Animal Diversity Web \citep{myer2000animal}, AnAge \citep{tacutu2018human}, and taxon-specific references \citep{ebert2021sharks, last2016rays}. When sources provided value ranges, the median was taken; when series of observations were available, the mean was used. Source priority followed: Animal Diversity Web $>$ IUCN $>$ AnAge. Validation used linear regression through the origin ($y \sim x$) at three feeding levels ($E(Y) = 0.5, 0.7, 0.9$), with 95\% CIs computed as regression coefficient $\pm$ 2$\times$ SE.

\subsection{Data and Code Availability}

The DEBBIES dataset (Version 5) is available on FigShare as CSV with metadata. MATLAB code packages are provided for cross-taxonomic analysis (computing all derived traits, $\lambda$, damping ratio, and elasticity for all species) and single-species analysis (computing the same quantities for one species across user-defined feeding levels).

\section{Conclusion}

DEBBIES provides a standardized resource for parameterizing DEB-IPMs across 185 ectotherm species in 18 orders, filling a critical gap in the demographic database landscape. By requiring only eight species-level life history traits, DEBBIES enables population modeling for data-deficient species excluded from conventional demographic analyses. Technical validation demonstrates strong agreement between DEB-IPM predictions and observed demographic quantities at ecologically realistic feeding levels. The dataset and code, freely available on FigShare, support cross-taxonomic life history comparisons, mechanistic population forecasting, and integration with Essential Biodiversity Variable frameworks for global biodiversity monitoring.

\bibliographystyle{plainnat}
\bibliography{references}

\end{document}
